\begin{problem}{Natural Gradient}{}{Modified from CS 229 Stanford 2018 - Problem set 3}
We consider the problem of fitting model parameters using gradient descent. Fitting model \\
parameters using Maximum Likelihood is equivalent to minimizing the KL divergence between the data and the model.
A nice property of KL divergence is that it is invariant to parametrization.
However, the gradient w.r.t the model parameters (i.e, direction of steepest descent)
is \emph{not invariant to parametrization}.

This is where the \textbf{natural gradient} comes into picture. It is best introduced in contrast with the
usual gradient descent. In the usual gradient descent, we first choose the direction by calculating
the gradient of the MLE objective w.r.t the parameters, and then move a magnitude of step size
(where size is measured in the parameter space) along that direction. Whereas in natural gradient,
we first choose a divergence amount by which we would like to move, in the \(D_{\mathrm{KL}} \) sense. This
effectively gives us a perimeter around the current parameters (of some arbitrary shape), such
that points along this perimeter correspond to distributions which are at an equal \(D_{\mathrm{KL}} \)-distance
away from the current parameter. Among the set of all distributions along this perimeter, we
move to the distribution that maximizes the objective (i.e minimize \(D_{\mathrm{KL}} \) between data and
itself) the most. This approach makes the optimization process invariant to parametrization.
That means, even if we chose a new arbitrary reparametrization, by starting from a particular
distribution, we always descend down the same sequence of distributions towards the optimum.
\begin{itemize}
    \item[(a)] Derive the natural gradient update rule.
    \item[(b)] Show natural gradient based optimization is actually equivalent to
          Newton’s method for Generalized Linear Models.
\end{itemize}

\end{problem}

\begin{solution}
    \begin{itemize}

        \item[(a)] The natural gradient update rule is:
              \[
                  d^* = \arg\max_d \ell(\theta + d) \quad \text{subject to} \quad D_{\text{KL}}(P_\theta || P_{\theta+d}) = c \quad (1)
              \]

              \begin{lemma}{Approximating $D_{\text{KL}}$ with Fisher Information}{fisher}
                  \[
                      D_{\text{KL}}(P_\theta || P_{\theta+d}) \approx \frac{1}{2} d^T \mathcal{I}(\theta) d
                  \]
              \end{lemma}

              Consider the Taylor approximation of $\ell(\theta + d)$ around $\theta$:
              \[
                  \ell(\theta + d) \approx \ell(\theta) + d^T \nabla_{\theta'} \ell(\theta') \Big|_{\theta'=\theta}
              \]

              Plugging this approximation into (1), we get the following optimization problem:
              \[
                  d^* = \arg\max_d \ell(\theta) + d^T \nabla_{\theta'} \ell(\theta') \Big|_{\theta'=\theta} \quad \text{s.t.} \quad \frac{1}{2} d^T \mathcal{I}(\theta) d = c
              \]

              We optimize the Lagrangian $\mathcal{L}(d, \lambda)$:
              \[
                  \mathcal{L}(d, \lambda) = \ell(\theta) + d^T \nabla_{\theta'} \ell(\theta') \Big|_{\theta'=\theta} - \lambda \left[ \frac{1}{2} d^T \mathcal{I}(\theta) d - c \right]
              \]

              \begin{itemize}
                  \item $\nabla_d \mathcal{L}(d, \lambda) = 0 \implies \nabla_{\theta'} \ell(\theta') \Big|_{\theta'=\theta} - \lambda \mathcal{I}(\theta) d = 0 \implies d = \frac{1}{\lambda} \mathcal{I}^{-1}(\theta) \nabla_{\theta'} \ell(\theta') \Big|_{\theta'=\theta}$
                  \item $\nabla_\lambda \mathcal{L}(d, \lambda) = 0 \implies \frac{1}{2} d^T \mathcal{I}(\theta) d = c \implies d^T \mathcal{I}(\theta) d = 2c$
              \end{itemize}

              \begin{align*}
                   & \implies d^T \nabla_{\theta'} \ell(\theta') \Big|_{\theta'=\theta} = 2c \implies \frac{1}{\lambda} \left( \nabla_{\theta'} \ell(\theta') \Big|_{\theta'=\theta} \right)^T \mathcal{I}^{-1}(\theta) \nabla_{\theta'} \ell(\theta') \Big|_{\theta'=\theta} = 2c                \\
                   & \implies \lambda = \frac{1}{\sqrt{2c}} \sqrt{\left( \nabla_{\theta'} \ell(\theta') \Big|_{\theta'=\theta} \right)^T \mathcal{I}^{-1}(\theta) \nabla_{\theta'} \ell(\theta') \Big|_{\theta'=\theta}}                                                                          \\
                   & \implies d^* = \frac{\sqrt{2c} \mathcal{I}^{-1}(\theta) \nabla_{\theta'} \ell(\theta') \Big|_{\theta'=\theta}}{\sqrt{\left( \nabla_{\theta'} \ell(\theta') \Big|_{\theta'=\theta} \right)^T \mathcal{I}^{-1}(\theta) \nabla_{\theta'} \ell(\theta') \Big|_{\theta'=\theta}}}
              \end{align*}

              \textbf{Proof of Lemma 1:}

              The KL divergence can be split into two expectations:
              \[
                  D_{\text{KL}}(P_\theta || P_{\theta+d}) = \underbrace{\mathbb{E}_{x \sim P_\theta(x)} [\log P_\theta(x)]}_{A} - \underbrace{\mathbb{E}_{x \sim P_\theta(x)} [\log P_{\theta+d}(x)]}_{B}
              \]

              Treating $\theta$ as a fixed point, $A$ is constant. We perform a second-order Taylor Expansion on $B$ around $\theta$. Let $\tilde{\theta} = \theta + d$:
              \begin{align*}
                  B & = \mathbb{E}_{x \sim P_\theta(x)} [\log P_{\tilde{\theta}}(x)]                                                                                                                                                                                                                                                                             \\
                    & \approx A + d^T \underbrace{\mathbb{E}_{x \sim P_\theta(x)} \left[ \nabla_{\theta'} \log P_{\theta'}(x) \Big|_{\theta'=\theta} \right]}_{\text{First-order term}} + \frac{1}{2} d^T \underbrace{\mathbb{E}_{x \sim P_\theta(x)} \left[ \nabla_{\theta'}^2 \log P_{\theta'}(x) \Big|_{\theta'=\theta} \right]}_{\text{Second-order term}} d
              \end{align*}

              \textbf{Step 1: Show the first-order term is zero.}
              \begin{align*}
                  \mathbb{E}_{x \sim P_\theta(x)} \left[ \nabla_{\theta'} \log P_{\theta'}(x) \Big|_{\theta'=\theta} \right] & = \mathbb{E}_{x \sim P_\theta(x)} \left[ \frac{\nabla_{\theta'} P_{\theta'}(x)}{P_{\theta'}(x)} \Bigg|_{\theta'=\theta} \right]         \\
                                                                                                                             & = \int_{-\infty}^{\infty} P_\theta(x) \frac{\nabla_{\theta'} P_{\theta'}(x) \Big|_{\theta'=\theta}}{P_\theta(x)} dx                     \\
                                                                                                                             & = \nabla_{\theta'} \left[ \int_{-\infty}^{\infty} P_{\theta'}(x) dx \right] \Bigg|_{\theta'=\theta} = \nabla_{\theta'}[1] = 0 \quad (*)
              \end{align*}

              \textbf{Step 2: Connect the second-order term to the Fisher Information Matrix $\mathcal{I}(\theta)$.}
              By definition, and using the result from Step 1 $(*)$:
              \begin{align*}
                  \mathcal{I}(\theta) & = \text{Cov}_{x \sim P_\theta(x)} \left[ \nabla_{\theta'} \log P_{\theta'}(x) \Big|_{\theta'=\theta} \right]                                                                  \\
                                      & = \mathbb{E}_{x \sim P_\theta(x)} \left[ \left(\nabla_{\theta'} \log P_{\theta'}(x)\right) \left(\nabla_{\theta'} \log P_{\theta'}(x)\right)^T \Big|_{\theta'=\theta} \right]
              \end{align*}

              Examining a single element $(i, j)$ of this matrix:
              \begin{align*}
                  [\mathcal{I}(\theta)]_{ij} & = \mathbb{E}_{x \sim P_\theta(x)} \left[ \left( \frac{1}{P_{\theta'}(x)} \right)^2 \frac{\partial P_{\theta'}(x)}{\partial \theta'_i} \frac{\partial P_{\theta'}(x)}{\partial \theta'_j} \Bigg|_{\theta'=\theta} \right] \\
                                             & = \int_{-\infty}^{\infty} P_\theta(x) \left[ \frac{1}{P_\theta(x)^2} \frac{\partial P_{\theta'}(x)}{\partial \theta'_i} \frac{\partial P_{\theta'}(x)}{\partial \theta'_j} \Bigg|_{\theta'=\theta} \right] dx \quad (**)
              \end{align*}

              We know that the integral of the second derivative of the probability density is zero:
              \[
                  \int_{-\infty}^{\infty} \frac{\partial^2 P_{\theta'}(x)}{\partial \theta'_i \partial \theta'_j} \Bigg|_{\theta'=\theta} dx = \frac{\partial^2}{\partial \theta'_i \partial \theta'_j} \left( \int_{-\infty}^{\infty} P_{\theta'}(x) dx \right) \Bigg|_{\theta'=\theta} = 0
              \]

              Subtracting this zero-valued integral from $(**)$ yields:
              \begin{align*}
                  [\mathcal{I}(\theta)]_{ij} & = \int_{-\infty}^{\infty} P_\theta(x) \left[ \frac{1}{P_\theta(x)^2} \frac{\partial P_{\theta'}(x)}{\partial \theta'_i} \frac{\partial P_{\theta'}(x)}{\partial \theta'_j} - \frac{1}{P_\theta(x)} \frac{\partial^2 P_{\theta'}(x)}{\partial \theta'_i \partial \theta'_j} \Bigg|_{\theta'=\theta} \right] dx \\
                                             & = \int_{-\infty}^{\infty} P_\theta(x) \left[ - \frac{\partial}{\partial \theta'_j} \left( \frac{1}{P_{\theta'}(x)} \frac{\partial P_{\theta'}(x)}{\partial \theta'_i} \right) \Bigg|_{\theta'=\theta} \right] dx                                                                                              \\
                                             & = \int_{-\infty}^{\infty} P_\theta(x) \left[ - \frac{\partial^2 \log P_{\theta'}(x)}{\partial \theta'_i \partial \theta'_j} \Bigg|_{\theta'=\theta} \right] dx                                                                                                                                                \\
                                             & = \mathbb{E}_{x \sim P_\theta(x)} \left[ - \nabla_{\theta'}^2 \log P_{\theta'}(x) \Big|_{\theta'=\theta} \right]_{ij}
              \end{align*}
              Thus, $\mathcal{I}(\theta) = - \mathbb{E}_{x \sim P_\theta(x)} \left[ \nabla_{\theta'}^2 \log P_{\theta'}(x) \Big|_{\theta'=\theta} \right]$.

              \textbf{Step 3: Final Substitution.}
              Substitute the results from Step 1 and Step 2 back into the Taylor expansion for $B$:
              \begin{align*}
                  B & \approx A + 0 + \frac{1}{2} d^T \mathbb{E}_{x \sim P_\theta(x)} \left[ \nabla_{\theta'}^2 \log P_{\theta'}(x) \Big|_{\theta'=\theta} \right] d \\
                    & = A - \frac{1}{2} d^T \mathcal{I}(\theta) d
              \end{align*}

              Finally, substitute $B$ back into the KL divergence equation:
              \[
                  D_{\text{KL}}(P_\theta || P_{\theta+d}) \approx A - \left( A - \frac{1}{2} d^T \mathcal{I}(\theta) d \right) = \frac{1}{2} d^T \mathcal{I}(\theta) d
              \]

        \item[(b)]
              We need to show that the direction of update of Newton's Method, and the direction of natural gradient, are exactly the same for GLMs.

              Update rule for Newton's method:
              \[
                  \theta := \theta - H^{-1} \nabla_\theta \ell(\theta)
              \]
              where $H_{ij} = \frac{\partial^2 \ell(\theta)}{\partial \theta_i \partial \theta_j}$ is the Empirical Hessian.

              Exponential family:
              \[
                  p(y; \eta) = b(y) \exp(\eta^T T(y) - a(\eta))
              \]

              For the sake of convenience, assume $\eta \in \mathbb{R}$ and $T(y) = y$. This makes the exponential family take the form:
              \[
                  p(y; \eta) = b(y) \exp(\eta y - a(\eta))
              \]

              Log likelihood on GLM models:
              \begin{align*}
                  \ell(\theta) & = \log \prod_{i=1}^m p(y^{(i)} | x^{(i)} ; \theta) = \sum_{i=1}^m \log p(y^{(i)} | x^{(i)} ; \theta) \\
                               & = \sum_{i=1}^m \left[ \log b(y^{(i)}) + \theta^T x^{(i)} y^{(i)} - a(\theta^T x^{(i)}) \right]
              \end{align*}

              Taking the first derivative:
              \begin{align*}
                  \implies \nabla_\theta \ell(\theta) & = \sum_{i=1}^m \nabla_\theta \left[ \log b(y^{(i)}) + \theta^T x^{(i)} y^{(i)} - a(\theta^T x^{(i)}) \right] \\
                                                      & = \sum_{i=1}^m x^{(i)} \left[ y^{(i)} - a'(\theta^T x^{(i)}) \right]
              \end{align*}

              Taking the second derivative:
              \[
                  \implies \nabla_\theta^2 \ell(\theta) = \sum_{i=1}^m -a''(\theta^T x^{(i)}) x^{(i)} (x^{(i)})^T
              \]

              $\implies \mathbb{E}_{y \sim p(y)}[H] = H$ since $H$ is constant wrt $y$.

              Now consider the update direction of the Natural gradient: $\mathcal{I}^{-1}(\theta) \nabla_\theta \ell(\theta)$.

              We have:
              \[
                  \mathcal{I}(\theta) = \mathbb{E}_{y \sim p(y)}[-H] = -H
              \]

              Compared to the Newton update, we see the update direction between the two methods matches for GLMs.
    \end{itemize}

\end{solution}